\scanpage{25}
\[
\Sigma(x)\subset R(x),\qquad\Sigma(x)\ne\varnothing.
\]
Cette notion ne dépend que de $C,R,C(j)\subset C\setminus C_0$.

\inserted{\uline{NB} On pose $\Sigma(x)=R(x)$ si $x\in\bar C(j)$. Donc on peut interpréter $\Sigma$ comme une application $C\setminus C_0\longrightarrow\mathcal P^*(C)$, $\Sigma(x)\subset R(x)$ et $\Sigma(x)\ne\varnothing$ $\forall x\in C\setminus C_0$, et $\Sigma(x)=R(x)$ si $x\in\bar C(j)$.}
\deleted{\uline{Stratégie \unsure{pour $j$} gagnante pour $x_0\in C$}}

Partie jouée en accord avec une Stratégie $\Sigma$ (ou \uline{compatible} avec $\Sigma$, ou $\Sigma$-partie) : cette notion est stable par composition. (On prend une sous-catégorie ayant mêmes objets, mais a priori moins de flèches)

\inserted{\uline{Lemme} Soit \deleted{\unsure{une stratégie}} $J'\subset J$, $\forall j\in J'$ soit $\Sigma_j$ une $j$-stratégie. Soit $x_0\in C$. Alors $\exists$ partie maximale d'origine $x_0$, comp. avec les $\Sigma_j$ $(j\in J')$.}

Stratégie $\Sigma$ (pour $j$) \uline{gagnante} pour $x_0\in C$ (ou \uline{position} $x_0$ \uline{gagnante} pour une Stratégie $\Sigma$) : toute partie maximale commençant en $x_0$, et jouée suivant $\Sigma$, est gagnée par $j$ (i.e. est finie, et son extrémité $x_n\in C_0$ est $\in G_0(j)$).
\margin{ne dépend que de $C,R,C(j),G_0(j)$}

\uline{Position \inserted{$x_0\in C$} gagnante pour} $j$ : telle qu'il existe une $j$-Stratégie qui soit gagnante pour $x_0$. Soit
\[
G(j)
\]
l'ens de ces positions. On a
\[
G(j)\cap C_0=G_0(j).
\]

[\uline{Proposition} Soit $J'\subset J$. Si $\bigcap_{j\in J'}G(j)\ne\varnothing$, alors $\bigcap_{j\in\unsure{J|J'}}G_0(j)\ne\varnothing$. En d'autres termes, si $\bigcap_{j\in J'}G_0(j)=\varnothing$, alors $\bigcap_{j\in J'}G(j)=\varnothing$ i.e. un $x_0\in C$ ne peut être gagnant pour tous les $j\in J'$.

Immédiat par le lemme.]
% archive foliation: 24
